\section{Метасвидетельства относительно решения}\label{sec:meta-evidence}
Сертификат корректного обучения не является сертификатом экономического выигрыша. Различим стоимость выбора контроллера, стоимость получения сведений о его стоимости и достаточность этих сведений для решения. Конечная метазадача сама является экземпляром модели исследования раздела~3; её сценарии --- пары задач с фиксированными размерами носителей.

\subsection{Выбор финансируемых исследовательских планов}
План задаёт список базовых схем правил, лимит обучения $\lambda$, резерв завершения $R$ и начальную плату $k$. Сначала выполняется оплачиваемый префикс обучения в пределах $\lambda$, затем текущий фронтир, неизрасходованный лимит и $R$ передаются статическому исполнителю с отсечением. Требуется доказательство, что передаваемый остаток не меньше статического потенциала обязательств. Запись префикса, включая любое ожидающее событие обучения, возвращается отдельно; ожидающая компиляция не является оставшимся обязательством по ответу. При нулевом лимите обучение отключено, включая плату за поиск в библиотеке.

\begin{theorem}[Финансируемые планы и метавыбор]\label{thm:meta-selection}
Каждый финансируемый план возвращает в точности каноническое членство для своей пары задач. Его фактическая модельная стоимость, включая начальную плату, не превосходит $k+\lambda+R$. Пусть для конечного каталога таких планов, индексированного действиями, таблица $U(a,w)$ ограничивает сверху фактическую стоимость завершения $C(a,w)$ каждого плана в каждом сценарии, и задана фиксированная плата контроллера $k_0$. Оптимизатор исследования раздела~3 минимизирует $k_0+W_E(t)$ при терминальной потере $U$; его прогнозируемая худшая стоимость ограничивает фактическое исполнение. Если всюду $U=C$, оптимизатор минимизирует фактическую худшую стоимость по допустимым метадеревьям.
\end{theorem}
\begin{proof}
Префикс обучения сохраняет корректность. Финансируемая передача завершается по теореме~\ref{thm:funded-search}; сохранение средств в обеих фазах даёт оценку и отделяет неизрасходованный резерв от затрат. На метауровне выбранное дерево исполняет только выбранный им план. В каждом мире его фактическая стоимость равна сумме реализованных плат за запросы и $C(a,w)$, ограниченной соответствующей суммой с $U$. Теорема~\ref{thm:investigation} оптимизирует эту верхнюю целевую функцию; точность таблицы превращает неравенство в фактическую оптимальность. См. \rc{perform\_complete}, \rc{actualBill\_le\_reservation}, \rc{optimizer\_bound\_optimal} и \rc{optimizer\_actual\_optimal} в \rc{MetaInvestigation.lean}.
\end{proof}
Таблица, доказательства финансирования и плата контроллера предоставляются, а не выводятся из бесплатного оракула. Явный свидетель уменьшает худшую модельную стоимость с $53$ до $25$ (включая объявленную плату метанаблюдения); офлайн-получение полной таблицы стоит $465$. При этих фиксированных платах строгая окупаемость требует не менее $17$ использований, поскольку $465+25n<53n$ равносильно $465<28n$. Это модельный свидетель, не прогноз будущей нагрузки. Необходимость получить таблицу стоимости мотивирует частичные свидетельства.

\subsection{Интервальные свидетельства и платное уточнение}
При фиксированных мирах, портах, ценах и глубине каждая неизвестная терминальная стоимость $C(a,w)$ представляется интервалом, индексированным доказательством:
\[
  L(a,w)\le C(a,w)\le U(a,w).
\]
Интервалы можно пересекать; нижние границы повышаются, верхние понижаются, фактическая стоимость остаётся покрытой. Пусть $V_L,V_C,V_U$ --- соответствующие минимаксные значения раздела~3, $W_C$ --- худшая стоимость при фактических потерях, а $t_U$ --- оптимизатор верхней модели.
\begin{theorem}[Достаточный интервальный сертификат]\label{thm:interval}
$V_L\le V_C\le V_U$, и для каждого допустимого конкурирующего дерева $t$
\[
 W_C(t_U)\le W_C(t)+(V_U-V_L).
\]
Зазор $V_U-V_L$ не может увеличиваться при уточнении интервалов. В частности, зазор не более $\epsilon$ достаточен для сертификации верхнего оптимизатора как $\epsilon$-близкого к оптимальному относительно фактической минимаксной цели.
\end{theorem}
\begin{proof}
Стоимости деревьев монотонны по терминальным потерям, а допустимость не меняется; поэтому соответствующие оптимальные значения упорядочены. Далее $W_C(t_U)\le V_U\le W_C(t)+V_U-V_L$, поскольку $V_L\le V_C\le W_C(t)$. Уточнение улучшает обе границы в соответствующих направлениях. Формальные утверждения: \rc{lower\_le\_true}, \rc{true\_le\_upper}, \rc{certified\_gap} и \rc{gap\_monotone} в \rc{MetaEvidenceAcquisition.lean}.
\end{proof}
Получение и оценка --- отдельные операции с положительными платами в предоставленном конечном расписании. Граница, индексированная фактической стоимостью, приобретается только после оплаты её котировки; исполнитель регистрирует платежи, монотонно обновляет свидетельства и останавливается при достаточном сертификате или исчерпании бюджета/расписания. Теорема \rc{assess\_complete} устанавливает полноту теста зазора, а не необходимое условие близости к оптимуму. Для конкретного семейства финансируемых планов приобретение шести из восьми точных ячеек сертифицирует значение $22$ при предоплате $2309$ против $3051$ за полное получение. Это свидетель расписания, установленного постфактум, а не оптимальный алгоритм приобретения. Предоплата по резерву не тождественна реализованной стоимости эксперимента.

\subsection{Достаточность для решения и оптимальное приобретение свидетельств}
Неопределённость стоимости не влечёт неопределённости решения. Например, в двух возможных моделях со стоимостями $(0,1)$ и $(100,101)$ для двух контроллеров первый оптимален в обеих, несмотря на большой разброс оптимальных значений. Поэтому используем конечный непустой каталог $\mathcal C$ допущенных контроллеров и список $E$ совместимых гипотез. Каждая гипотеза $h$ задаёт стоимость $J_h(c)$ и эталон $v_h\le J_h(c)$ для всех $c\in\mathcal C$. Эталон может быть достигаемым минимумом каталога, но не обязан быть оптимумом по более широкому классу. Определим
\[
 r_h(c)=J_h(c)-v_h,\qquad R(E)=\min_{c\in\mathcal C}\max_{h\in E}r_h(c),
\]
с максимумом пустого списка, равным нулю.
\begin{theorem}[Достаточность относительно решения]\label{thm:decision}
$R(E)\le\epsilon$ тогда и только тогда, когда некоторый допущенный контроллер удовлетворяет $J_h(c)\le v_h+\epsilon$ для каждого $h\in E$. Для эталона каталога это эквивалентно $J_h(c)\le J_h(c')+\epsilon$ для всех $h\in E,c'\in\mathcal C$. Уточнение списка гипотез не увеличивает $R$.
\end{theorem}
\begin{proof}
Минимум достигается конечным выбором, а ограничение максимума эквивалентно поточечным неравенствам. Для каталога в одну сторону используется нижняя граница минимума каталога на всех конкурентах, в обратную --- достигающий минимум конкурент для каждой гипотезы. Уточнение подмножеством может лишь уменьшать худшее сожаление каждого контроллера, а значит, и минимум. Это \rc{residual\_iff}, \rc{catalogue\_robust\_iff} и \rc{residual\_subset} в \rc{DecisionEvidence.lean}.
\end{proof}
Попарные границы стоимости могут сертифицировать решение без идентификации модели стоимости. Принадлежность действительной гипотезы --- явная внешняя предпосылка: детерминированный ответ порта $q$ фильтрует $E$, сохраняя действительную гипотезу, если она присутствовала изначально (\rc{refine\_truth}). Используя терминальную потерю $r_h(c)$ в модели исследования, получаем оптимальное дерево приобретения по сумме затрат на запросы и сожаления в худшем случае, с терминальным значением $R(E)$ и рекурсией Беллмана раздела~3. Это \rc{acquisition\_optimal} и \rc{acquisition\_bellman}. Фактическая полная стоимость при внешней подготовке $k$ в точности равна $k+v_h+$ реализованная сумма платы за запросы и сожаления; минимизация худшего сожаления в общем случае не минимизирует худшую фактическую стоимость, если $v_h$ меняется (\rc{total\_bill\_identity}).

Жёсткий контракт отличается от мягкой суммы. Используя ограниченную конструкцию теоремы~\ref{thm:investigation}, \rc{hardAssess} успешен тогда и только тогда, когда существует допустимое дерево приобретения с расходами на наблюдения не более $b$ и терминальным сожалением не более $\epsilon$ для каждой совместимой гипотезы. При достижимости его дерево также минимизирует худшие расходы на наблюдения среди допустимых деревьев, удовлетворяющих допуску сожаления (\rc{hard\_assess\_complete}, \rc{hard\_minimal\_observation}). Эти гарантии требуют принадлежности действительной гипотезы и заданного конечного каталога, а не автоматической генерации гипотез или бесплатной валидации. В частности, предварительный выбор контроллера не разрешает его установку в частично исследованный текущий фронтир. Это отдельная задача замены раздела~6.
